\section{Test Case 9: 1,000-bead CPU/GPU performance benchmark}

Test Case 9 is a performance-oriented three-dimensional Maxwell simulation
derived from Test Case 7. It initializes 1,000 beads uniformly over the full
16 cm distance from the nozzle to the collector and executes 1,000 integration
steps. Insertion, removal, multiple-step Coulomb summation, and dynamic
refinement are disabled, so the workload remains fixed.

The material, aerodynamic, gravity, and perturbation parameters are retained
from Test Case 7. The orthogonal rotating field is replaced by a constant
axial electric potential. Evaporation is disabled because its direct Coulomb
routine is not yet accelerated by the initial OpenACC implementation. The
case therefore exercises the direct $O(N^2)$ Coulomb kernel without an
evaporation fallback or topology changes.

Test Case 9 is intended for performance comparison rather than as an
additional physical validation baseline. CPU and accelerator runs should be
compared using the terminal value \texttt{Time-integration loop wall time},
which excludes input processing, initial allocation, and accelerator
initialization. Its input is stored in
\texttt{examples/input-9/input.dat}.

After compiling separate CPU and accelerator executables, copy each executable
and \texttt{examples/input-9/input.dat} into a separate run directory. Execute
the two cases sequentially on an otherwise idle system and compare their
\texttt{Time-integration loop wall time} values. The compiler version, GPU
model, GPU compute capability, and build options should accompany every
reported benchmark.

An initial reference measurement with NVFORTRAN 24.3 and one NVIDIA A30
reported 40.659856 s for the CPU executable and 3.785055 s for the OpenACC
executable, corresponding to a speedup of 10.742. Both runs retained exactly
1,000 beads and agreed within the numerical regression tolerances. These
timings are indicative and should be remeasured on each target system.

An optional internal subroutine profiler is enabled at run time with
\texttt{JETSPIN\_PROFILE=1 ./main.x}; no rebuild is required. The report
separates the integrator, Coulomb force, bead topology operations, statistics,
scheduled output, and restart output. The Coulomb region is nested within the
integrator total and the two values must not be added.

The report is written only to the terminal and gives cumulative wall time,
percentage of the complete temporal loop, and call count for each explicitly
instrumented region. It is not an automatic profile of every Fortran
subroutine. The fourth-order Runge--Kutta integration used by this test makes
four Coulomb evaluations per step, yielding 4,000 calls. Clock reads add a
small overhead, so profiling must have the same enabled or disabled state in
both runs being compared. In MPI execution, the detailed table is local to
rank zero rather than a maximum over ranks; the synchronized complete-loop
timer remains the appropriate MPI scaling measurement.

In an initial profiled run, the NVFORTRAN CPU loop required 39.813621 s,
of which 38.232985 s (96.03\%) was spent in 4,000 Coulomb evaluations.
The A30 loop required 3.320607 s, with 2.164045 s (65.17\%) in the same
Coulomb region. The corresponding speedups were approximately 11.99 for the
complete loop and 17.67 for Coulomb alone. Repeating this profile after each
porting milestone provides a direct view of time moving from the Coulomb
kernel to remaining host work, device transfers, or other force terms.

The complete sampled CPU and A30 OpenACC numerical outputs from this initial
measurement are versioned in
\texttt{tests/performance/test9/}. The accompanying comparison script checks
later CPU and GPU results with relative tolerance $10^{-6}$ and absolute
tolerance $10^{-9}$. These records remain outside the automatic regression
matrix because of the benchmark's longer execution time.
